% Ad Atticum 9.10 (ad-atticum-09-10) --- English translation
% Translated by Claude (Anthropic) from the Latin (Perseus).
% Style guide: STYLE.md.

\ciceroLetterOpener{CICERO ATTICO salutem}

\ciceroSection{1} I had nothing to write. For neither had I heard anything new, and to all of yours I had written back the day before. But since distress was depriving me not only of sleep, but would not even let me lie awake without the highest pain, I set myself to write to you this something or other, with no theme proposed, as if to talk with you --- which is the one thing in which I find rest.

\ciceroSection{2} I seem to myself to have been out of my mind from the start, and this one thing tortures me --- that, when Pompey was slipping in every matter, or rather hurtling down, I did not follow him like one common soldier in the ranks. I saw the man on the fourteenth day before the Kalends of February: he was full of dread. That very day I felt what he was about. Never afterward did he please me, nor did he ever stop committing one fault after another. To me, all this while, he wrote nothing; he thought of nothing but flight. What more is there? As, in love-affairs \textit{[Greek: en tois er\=otikois]}, things filthy, insipid, indecorous estrange the lover, so the ugliness of his flight and carelessness turned me away from love. For he did nothing worthy that I should join myself as the companion of his flight. Now the love comes back up; now I cannot bear the longing; now no books, no letters, no learning is of use to me. Day and night, like that bird, I gaze out upon the sea and long to fly away. I am paying, I am paying the penalty for my rashness. And yet what rashness was it? what have I done that was not most carefully weighed? For if nothing but flight had been the object, I should have fled most willingly; but the kind of war --- the most cruel and the greatest, the like of which men do not yet see what it will be --- this I shuddered at. What threats against the towns, what against good men by name, what, in short, against all who had remained! How often the line: ``Sulla could; cannot I?''

\ciceroSection{3} But these things have stuck with me. Tarquin acted ill who led Porsena and Octavius Mamilius against the fatherland; impiously, Coriolanus, who sought aid from the Volscians; rightly, Themistocles, who preferred to die; criminally, Hippias the son of Pisistratus, who fell in the battle of Marathon bearing arms against his fatherland; while Sulla, Marius, Cinna acted rightly, indeed perhaps by right --- but what was more cruel than their victory, what more deadly? This kind of war I fled, and the more so because I saw even crueler things being thought of and prepared. Was I, whom certain men have called the preserver of this city, whom they have called its parent, to lead against it the forces of the Getae and the Armenians and the Colchians? Was I to bring famine, devastation, upon my fellow citizens, upon Italy? I considered, first, that this man is mortal, and then, that he can be quenched in many ways; but our city and our people I thought we ought to save for immortality, so far as it lay in us. And still a certain hope sustained me, that some agreement would be reached rather than that this man should commit so much wickedness or that man so much disgrace. Now the whole matter is altered, my whole mind altered. The sun, as it is said in some letter of yours, seems to me to have fallen out of the world. As they say of one who is sick that, while there is breath, there is hope, so I, as long as Pompey was in Italy, did not stop hoping. These things, these things, have failed me; and, to speak the truth, my age, now bending down from long labours to leisure, has been softened by the pleasure of my domestic affairs. Now, if the experiment must be made even at peril, I shall certainly make it, so that I may fly away from here. Perhaps I ought to have done it earlier; but the things you wrote held me back, and above all your authority.

\ciceroSection{4} For when I had come to this point, I unrolled the volume of your letters which I keep under seal and guard most carefully. There was, then, in the one you sent on the tenth day before the Kalends of February, this passage: ``But let us see what Gnaeus is doing, and where his calculations are tending. If that man leaves Italy, he will be acting wholly badly and, as I judge, without reckoning \textit{[Greek: alogist\=os]}; but then, and only then, will our counsels have to be changed.'' This you write on the fourth day after we left the City. Then on the eighth day before the Kalends of February: ``Only let our Gnaeus not leave Italy as he left the City --- without reckoning \textit{[Greek: alogist\=os]}.'' On the same day you give another letter, in which you answer me, who was consulting you, very plainly. The wording is as follows. ``But I come to your question. If Gnaeus quits Italy, I think you must return to the City; for what end can there be to a wandering abroad?'' This stuck plainly with me; and now too I see an unending war joined to the most wretched flight, which you call by the soft name of a ``foreign tour'' \textit{[Greek: hypokoriz\=ei]}.

\ciceroSection{5} Then follows the oracle \textit{[Greek: chr\=esmos]} of the sixth day before the Kalends of February: ``I, if Pompey remains in Italy and the matter does not come to an agreement, think the war will be longer; but if he leaves Italy, I judge that a war without truce \textit{[Greek: aspondon]} is being prepared for the next generation.'' Am I then to be made a sharer, an ally, a helper in this war, when it is both without truce \textit{[Greek: aspondon]} and against fellow citizens? Then on the seventh day before the Ides of February, when you were now hearing more about Pompey's plans, you close a certain letter in this way: ``I, indeed, would not be your authority, if Pompey leaves Italy, that you flee too. For you will be acting at the greatest risk, and you will be of no use to the commonwealth; whereas you will be able to be of use to it later, if you remain.'' What patriot \textit{[Greek: philopatrin]} and statesman \textit{[Greek: politikon]} would not be moved by the authority of a prudent man and a friend admonishing him so?

\ciceroSection{6} Next, on the third day before the Ides of February, you answer me again, as I was consulting you, thus: ``As to what you ask of me, whether I would defend flight or think delay the more useful: I, for my present part, think a sudden departure and headlong setting-off pernicious and dangerous both for you and for Gnaeus himself, and consider that it is better that you be apart and on watch; but, by Hercules, I think it is shameful for us to be thinking of flight.'' This shameful thing our Gnaeus has been thinking of for two years past. Thus his mind has long now been yearning-to-be-a-Sulla, yearning-to-proscribe. Then, I take it, when you had written something to me in more general terms \textit{[Greek: genik\=oteron]}, and I had thought you were signifying something to me about my quitting Italy, you carefully repudiate this on the eleventh day before the Kalends of March: ``I have not, by any letter, signified that, if Gnaeus leaves Italy, you should leave with him; or, if I did so signify --- I say not that I was inconsistent, but that I was mad.'' In the same letter, in another place: ``Nothing is left but flight; of which I do not at all think you ought to be a sharer, nor have I ever thought it.''

\ciceroSection{7} The whole of this deliberation, however, you unroll more carefully in the letter sent on the eighth day before the Kalends of March: ``If Manius Lepidus and Lucius Volcacius remain, I think you must remain; in such a way that, if Pompey is safe and has come to a stand somewhere, you should leave this Necyia \textit{[Greek: nekuian]} and let yourself be beaten in the contest with him more easily than reign with this other man in that confusion which is plainly going to come.'' You argue many things consistent with this opinion. Then at the end: ``What if,'' you say, ``Lepidus and Volcacius depart? Plainly, I am at a loss \textit{[Greek: apor\=o]}. So whatever turns out, and whatever you do, I shall think must be accepted \textit{[Greek: sterkteon]}.'' If you doubted then, certainly now you do not doubt, with those men remaining.

\ciceroSection{8} Then in the flight itself, on the fifth day before the Kalends of March: ``In the meantime, I have no doubt that you will be staying at the Formian villa. There you will wait out what is to come \textit{[Greek: to mellon karadok\=eseis]} most conveniently.'' And on the Kalends of March, when he had been at Brundisium for the fifth day now, ``Then we shall be able to deliberate --- with the matter, of course, not whole, but certainly less broken than if you had flung yourself in alongside him.'' Then on the fourth day before the Nones of March, when you were writing briefly because you were at the very moment of seizure \textit{[Greek: hypo t\=en l\=epsin]}, still you put this: ``Tomorrow I shall write more, and to all your points; this, however, I shall say --- that I do not regret my advice about your staying. Although I am in great anxiety, still, because I think there is less evil in it than in that setting-off, I keep my opinion, and rejoice that you have stayed.''

\ciceroSection{9} And when I was already in real anguish, and afraid that some disgrace had been admitted on my part, on the third day before the Nones of March: ``Still, I do not bear it ill that you are not with Pompey. Afterwards, if it shall be needful, it will not be difficult; and to him, at whatever time it comes about, it will be welcome \textit{[Greek: asmeniston]}. But this I say with the proviso that, if this man here carries out everything in the manner he began --- sincerely, with moderation, with prudence --- then I shall examine the matter strictly, and shall provide for our advantage with more deliberate care.''

\ciceroSection{10} On the seventh day before the Ides of March you write that our Peducaeus too approves my having kept quiet --- and his authority weighs much with me. With these writings of yours I console myself, in such a way that I think no fault has so far been committed by me. Only, do you defend your authority: against me there is no need, but I am in want of other witnesses. If I have done no wrong, I shall defend what remains. Encourage me to it, and help me throughout with your thinking. Here nothing was yet being heard about Caesar's return. By this letter, however, I have achieved this much: I have read through all yours, and have found rest in that.
